\subsection{Morphisms of commutative Banach algebras}

PROPOSITION 9. --- Let A be a Banach algebra, and let B be a commutative semisimple Banach algebra. Every morphism from the underlying algebra of A into the underlying algebra of B is continuous.

Let \(h \colon \mathrm{A} \to \mathrm{B}\) an algebra homomorphism and let \((a, b) \in \mathrm{A} \times \mathrm{B}\) a point adherent to the graph \(\Gamma\) of \(h\) . Let \(\chi \in \mathsf{X}'(\mathrm{B})\) . The function \(x \mapsto \chi(h(x))\) of \(\mathrm{A}\) into \(\mathbf{C}\) is an algebra homomorphism, hence is continuous (I, p.~29, th. 1). The map from \(\mathrm{A} \times \mathrm{B}\) into \(\mathbf{C}\) given by \((x, y) \mapsto \chi(h(x)) - \chi(y)\) is then continuous, it vanishes on \(\Gamma\) therefore zero at \((a, b)\) We thus have \(\chi(h(a)) = \chi(b)\) for all \(\chi \in \mathsf{X}'(\mathrm{B})\) . Since B is semisimple, we have \(h(a) = b\) . Thus the graph of \(h\) is closed and therefore \(h\) is continuous (EVT, I, p.~19, cor. 5).

COROLLARY. --- On a commutative complex algebra without radical, two norms defining Banach algebra structures are equivalent.

It suffices to apply Prop. 9 to the identity map of the algebra.

Let A and B be commutative Banach algebras. By No. 7 of I, p.~9, if \(h\colon A\to B\) is a surjective morphism, \(\mathsf{X}'(h)\) is a homeomorphism from \(\mathsf{X}'(\mathsf{B})\) on a closed subspace of \(\mathsf{X}'(\mathsf{A})\) which maps 0 to 0. (In the case where \(h\) is the injection of a subalgebra into \(\mathcal{C}(\mathsf{X})\) , the map \(\mathsf{X}'(h)\) is injective under much weaker hypotheses, cf.~I, p.~142, prop. 1, d).)

Now let \(h \colon \mathrm{A} \to \mathrm{B}\) an injective morphism. In general, \(\mathsf{X}'(h)\) is not surjective, but the following proposition provides a necessary condition for this to be the case.

PROPOSITION 10. --- Let A and B be unital commutative Banach algebras, \(h\colon A\to B\) a unital algebra homomorphism, not necessarily continuous. If \(\mathsf{X}(h)\) is surjective, then \(h(\mathrm{A})\) is a full subalgebra of B.

Let \(x \in A\) such that \(h(x)\) be invertible in B. For every \(\chi \in \mathsf{X}(\mathsf{A})\) , there exists \(\xi \in \mathsf{X}(\mathsf{B})\) such that \(\chi = \mathsf{X}(h)(\xi)\) , hence \(\chi(x) = \xi(h(x)) \neq 0\) . Prop. 6 of I, p.~37 then shows that x is invertible in A, and therefore that \(h(x)\) is invertible in \(h(\mathsf{A})\) .

The necessary condition of the proposition is not sufficient, even if h is isometric (I, p.~168, exerc. 14). We nevertheless have the following result:

PROPOSITION 11. --- Let A and B be commutative unital Banach algebras, let a be an element of A and let h: A → B be an injective unital morphism (not necessarily continuous). Suppose that the closed full subalgebra of A generated by a is equal to A. The following conditions are equivalent:

\begin{enumerate}
\def\labelenumi{(\roman{enumi})}
\item
  \(\mathsf{X}(h)\) is surjective;
\item
  \(h(\mathrm{A})\) is a full subalgebra of B;
\item
  \(\mathrm{Sp}_{\mathrm{A}}(a)=\mathrm{Sp}_{\mathrm{B}}(h(a)).\)
\item
  \(\Longrightarrow\) (ii) follows from prop. 10.
\item
  \(\Longrightarrow\) (iii) follows from the formula \(\mathrm{Sp}_{\mathrm{A}}(a) = \mathrm{Sp}_{h(\mathrm{A})}(h(a)) = \mathrm{Sp}_{\mathrm{B}}(h(a))\) , valid since \(h(\mathrm{A})\) is a full subalgebra of B.
\item
  \(\Longrightarrow\) (i) by formula (4) of no. 6 of I, p.~6, we have the commutative diagram
\end{enumerate}

\[
\begin{array}{c} \mathsf {X} (\mathrm{B}) \xrightarrow {\mathsf {X} (h)} \mathsf {X} (\mathrm{A}) \\ \Big \downarrow^ {\mathcal {G} _ {\mathrm{B}} (h (a))} \qquad \qquad \qquad \Big \downarrow^ {\mathcal {G} _ {\mathrm{A}} (a)} \\ \operatorname{Sp} _ {\mathrm{B}} (h (a)) \xrightarrow {i} \operatorname{Sp} _ {\mathrm{A}} (a) \end{array}
\]

where the vertical arrows denote surjective maps (I, p.~37, prop. 6) and \(i\) is the canonical inclusion. Hypothesis (iii) means that \(i\) is bijective. Moreover, the surjective map \(\mathcal{G}_{\mathrm{A}}(a)\colon \mathsf{X}(\mathrm{A})\to \mathrm{Sp}_{\mathrm{A}}(a)\) is bijective: indeed, for every characters \(\chi_{1}\) and \(\chi_{2}\) of A, the set \(\mathrm{A}_{\chi_1,\chi_2} = \{x\in \mathrm{A}\mid \chi_1(x) = \chi_2(x)\}\) is a closed full subalgebra of A. By hypothesis, we therefore have \(\mathrm{A}_{\chi_1,\chi_2} = \mathrm{A}\) if \(a\in \mathrm{A}_{\chi_1,\chi_2}\) , that is, \(\chi_{1} = \chi_{2}\) if \(\mathcal{G}_{\mathrm{A}}(a)\chi_{1} = \mathcal{G}_{\mathrm{A}}(a)\chi_{2}\) . The diagram then implies that the map \(\mathsf{X}(h)\) is surjective.

